\section{评价横截性}
\label{sec:evaluation-transversality}
\stfhSubIndexRangeStart{I290}{heng^jie^xing}{横截性（transversality）}{ping^jia^heng^jie^xing}{评价横截性（evaluation transversality）}

本节给出伪全纯曲线应用中的一个重要工具，即评价映射横截性。这种横截性在计数不变量的定义中起关键作用，例如格罗莫夫--威腾不变量的构造以及伪全纯曲线的粘合分析。

严格证明这种评价横截性颇为微妙。Lê Hồng Vân 与小野薫 (LO96)，以及 Ke Zhu 与吴龙根 (OhZ09)，在研究一阶 jet 评价横截性时采用了一种方法；我们把它用于这里更简单的零阶 jet 评价横截性，并给出概念上典范的证明。这个证明建立在点支集分布的标准结构定理之上。

(MSa04) 给出了另一种更具几何意味的方法。我们的证明具有自然性，因而有一项优势：除去使用 $J$-全纯曲线的\emph{全纯 jet} 之外，同一方案基本上可以推广到高阶 jet 横截性。该推广见 (Oh11a)。

先回忆第 4.5 节 (GS68) 中的点支集分布结构定理，即 (Rud73) 的定理 6.25；其证明请读者参阅原文。

\begin{theorem}[点支集分布]\label{thm:distribution-point-support}
\stfhIndex{I103}{dian^zhi^ji^fen^bu}{点支集分布（distribution with point support）}
设 $\psi$ 是开集 $\Omega\subset\mathbb R^n$ 上的分布，满足
$\operatorname{supp}\psi=\{p\}$ 且具有有限阶 $N<\infty$。则
\[
 \psi=\sum_{|\alpha|\leq N}D^\alpha\delta_p,
\]
其中 $\delta_p$ 表示 $p$ 处的狄拉克--delta 函数，而
$\alpha=(\alpha_1,\ldots,\alpha_n)$ 表示多重指标。
\end{theorem}

在下面的讨论中，为免同坐标 $z$ 混淆，把此前一直记作
$\{z_1,\ldots,z_m\}$ 的标记点集记成粗体
$\boldsymbol z$。考虑模空间
\[
 \widetilde{\mathcal M}_{g,m}(J;\beta)=
 \{((j,u),\boldsymbol z)\mid u:\Sigma\to M,
 \ \bar\partial_{(j,J)}u=0,
 \ \boldsymbol z=(z_1,\ldots,z_m)\in
 \widetilde{\operatorname{Conf}}_m(\Sigma)\}.
\]
评价映射
$\operatorname{ev}_i:\widetilde{\mathcal M}_{g,m}(J;\beta)\to M$ 定义为
\[
 \operatorname{ev}_i((j,u),\boldsymbol z)=u(z_i).
\]
再定义普适模空间
\[
 \widetilde{\mathcal M}_{g,m}(\beta)=
 \bigcup_{J\in\mathcal J_\omega}
 \widetilde{\mathcal M}_{g,m}(J;\beta)\longrightarrow\mathcal J_\omega,
\]
并以 $\widetilde{\mathcal M}_{g,m}^{\mathrm{inj}}(\beta)$ 表示其中由在某处单射曲线组成的开子集。
下面是本节的主定理。

\begin{theorem}[评价横截性]\label{thm:evaluation-transversality}
在每个
$((j,u),J,\boldsymbol z)\in
\widetilde{\mathcal M}_{g,m}^{\mathrm{inj}}(\beta)$ 处，对所有
$i=1,\ldots,m$，导数
\[
 D(\operatorname{ev}_i):
 T_{((j,u),J,\boldsymbol z)}\widetilde{\mathcal M}_{g,m}(\beta)
 \longrightarrow T_{u(z_i)}M
\]
都是满射。
\end{theorem}

考虑映射
\[
 \Upsilon:\mathcal J_\omega\times\mathcal F_{g,\ell}(\Sigma,M;\beta)
 \longrightarrow\mathcal H''\times M^\ell,
 \qquad
 (J,(j,u),\boldsymbol z)\longmapsto
 (\bar\partial_{(j,J)}u,\{u(z_i)\}).
\]
对 $J\in\mathcal J_\omega$，标准模空间 $\mathcal M_\ell(M,J;\beta)$ 的并正是
\begin{equation}
 \Upsilon^{-1}(o_{\mathcal H''}\times M^\ell)/\operatorname{Aut}(\Sigma),
 \label{eq:standard-moduli-as-zero-set}
\end{equation}
其中 $o_{\mathcal H''}$ 是上述丛 $\mathcal H''$ 的零截面，而
$\operatorname{Aut}(\Sigma)$ 对任意 $j$ 都通过共形重参数化作用于
$(\Sigma,j),u$。

先考虑一个标记点的情形，并把相应映射记作 $\Upsilon_1$。有
\[
 \widetilde{\mathcal M}_{g,1}(M;\beta)
 =\Upsilon_1^{-1}(o_{\mathcal H''}\times M),
 \qquad
 \widetilde{\mathcal M}_{g,1}(M,J;\beta)
 =\widetilde{\mathcal M}_1(M;\beta)\cap\pi_2^{-1}(J).
\]

\begin{proposition}\label{prop:one-point-evaluation-transverse}
对任意给定点 $p\in M$，映射 $\Upsilon_1$ 横截于子流形
\[
 o_{\mathcal H''}\times\{p\}\subset\mathcal H''\times M.
\]
\end{proposition}

\begin{proof}
它的线性化
$D\Upsilon_1(J,(j,u),z)$ 是映射
\begin{equation}
 (B,(b,\xi),v)\longmapsto
 \bigl(D_{J,(j,u)}\bar\partial(B,(b,\xi)),
 \ \xi(u(z))+du(z)(v)\bigr),
 \label{eq:linearization-one-point-evaluation}
\end{equation}
其中
$B\in T_J\mathcal J_\omega$、$b\in T_j\mathcal M(\Sigma)$、
$v\in T_z\Sigma$ 且 $\xi\in T_u\mathcal F(\Sigma,M;\beta)$。在
$W^{1,p}$ 上，这定义线性映射
\[
 T_J\mathcal J_\omega\times T_j\mathcal M(\Sigma)
 \times T_u\mathcal F(\Sigma,M;\beta)\times T_z\Sigma
 \longrightarrow
 \Omega_{(j,J)}^{(0,1)}(u^*TM)\times T_{u(z)}M.
\]
不过，为使 $\Upsilon_1$ 可微，需要至少取 $k\geq2$，用
$W^{k,p}$ 范数完备化 $\mathcal F(\Sigma,M;\beta)$。

我们要证明这个线性映射在每个
$(u,z_0)\in\widetilde{\mathcal M}_1(M;J,\beta)$ 处都是满射；也就是在满足
\[
 \bar\partial_{(j,J)}u=0,
 \qquad u(z_0)=x
\]
的每对 $(u,z_0)$ 处都满射。为此，对给定的
$\gamma\in\Omega_{(j,J)}^{(0,1)}(u^*TM)$ 与
$X_0\in T_{u(z_0)}M$，必须研究方程组
\begin{equation}
 D_{J,(j,u)}\bar\partial(B,(b,\xi))=\gamma,
 \qquad \xi(u(z_0))+du(v)=X_0
 \label{eq:evaluation-surjectivity-system}
\end{equation}
的可解性。

在评价横截性的研究中，定义域复结构 $j$ 并不起多大作用。除它作为参数出现之外，我们的计算始终用不到它。因此在整个证明中固定 $j$，只需考虑
$b=0=v$ 的情形。上述方程便化为
\begin{equation}
 D_{J,u}\bar\partial(B,\xi)=\gamma,
 \qquad \xi(u(z_0))=X_0.
 \label{eq:reduced-evaluation-system}
\end{equation}

现在对 $\xi\in W^{2,p}$ 研究\eqref{eq:reduced-evaluation-system}。把
\[
 \Omega_{1,p}^{(0,1)}(u^*TM)\times T_{u(z_0)}M=:\mathcal B_0
\]
看作赋予范数 $\|\cdot\|_{1,p}+|\cdot|$ 的巴拿赫空间，其中
$|\cdot|$ 是由 $T_xM\cong\mathbb C^n$ 上某个内积诱导的任意范数。我们将证明
\eqref{eq:linearization-one-point-evaluation} 限制在形如
$(B,(0,\xi),0)$ 的元素上后，作为映射
\[
 T_J\mathcal J_\omega\times\Omega_{2,p}^0(u^*TM)
 \longrightarrow
 \Omega_{1,p}^{(0,1)}(u^*TM)\times T_{u(z_0)}M
\]
是满射；这里 $(u,j,z_0,J)\in
\Upsilon_1^{-1}(o_{\mathcal H''}\times\{p\})$。
为澄清记号，用 $\langle\cdot,\cdot\rangle$ 表示自然配对
\[
 \Omega_{1,p}^{(0,1)}(u^*TM)\times
 \bigl(\Omega_{1,p}^{(0,1)}(u^*TM)\bigr)^*\longrightarrow\mathbb R,
\]
用 $(\cdot,\cdot)_x$ 表示 $T_xM$ 上的内积。

先证明像在 $\mathcal B_0$ 中稠密。设
\[
 (\eta,Y_p)\in
 \bigl(\Omega_{1,p}^{(0,1)}(u^*TM)\bigr)^*\times T_pM
\]
满足
\begin{equation}
 \left\langle D_u\bar\partial_{(j,J)}\xi+
 \frac12B\cdot du\circ j,\eta\right\rangle
 +\langle\xi,\delta_{z_0}Y_p\rangle=0
 \label{eq:annihilator-evaluation-linearization}
\end{equation}
对所有 $\xi\in\Omega_{2,p}^0(u^*TM)$ 与 $B$ 成立；其中
$\delta_{z_0}$ 是支于 $z_0$ 的狄拉克--delta 函数。不失一般性，可以像此前一样假设
$\xi$ 光滑，因为
$C^\infty(u^*TM)\hookrightarrow\Omega_{2,p}^0(u^*TM)$ 稠密。我们要证明
$\eta=0=Y_p$。

在\eqref{eq:annihilator-evaluation-linearization}中令 $B=0$，得到
\begin{equation}
 \langle D_u\bar\partial_{(j,J)}\xi,\eta\rangle
 +\langle\xi,\delta_{z_0}Y_p\rangle=0
 \quad\text{对所有 }\xi\in C^\infty,
 \label{eq:distribution-adjoint-with-delta}
\end{equation}
所以按分布导数的定义，$\eta$ 满足
\[
 (D_u\bar\partial_{(j,J)})^\dagger\eta-\delta_{z_0}Y_p=0,
\]
即
\[
 (D_u\bar\partial_{(j,J)})^\dagger\eta=\delta_{z_0}Y_p.
\]
形式伴随 $(D_u\bar\partial_{(j,J)})^\dagger$ 与
$D_u\bar\partial_{(j,J)}$ 有相同符号，因而是一阶椭圆微分算子。由椭圆正则性，
$\eta$ 在 $\Sigma\setminus\{z_0\}$ 上是经典解。还回忆
$\bar\partial^\dagger=-\partial$。

另一方面，在\eqref{eq:annihilator-evaluation-linearization}中令 $\xi=0$，得到
\begin{equation}
 \langle B\cdot du\circ j,\eta\rangle=0
 \label{eq:annihilator-almost-complex-evaluation}
\end{equation}
对所有 $B\in T_J\mathcal J_\omega$ 成立。前一节横截性证明所用的论证表明，在
$\Sigma\setminus\{z_0\}$ 的任意在某处单射点的小邻域中 $\eta=0$。由于假设
$u$ 在某处单射，而且在给定假设下其在某处单射点集在定义域中既开又稠密，这样的点存在。应用唯一延拓定理可得
$\eta=0$ 于 $\Sigma\setminus\{z_0\}$，因而 $\eta$ 作为 $\Sigma$ 上的分布，其支集包含在单点集
$\{z_0\}$ 中。

下面的引理将结束证明。

\begin{lemma}\label{lem:distributional-adjoint-continuous}
$\eta$ 是 $(D_u\bar\partial_{(j,J)})^\dagger\eta=0$ 在
$\Sigma$ 上的分布解，因而连续。特别地，在
$\bigl(\Omega_{1,p}^{(0,1)}(u^*TM)\bigr)^*$ 中有 $\eta=0$。
\end{lemma}

一旦知道 $\eta=0$，方程\eqref{eq:annihilator-evaluation-linearization}便化为有限维方程
\begin{equation}
 (\xi(z_0),Y_p)_{z_0}=0.
 \label{eq:finite-dimensional-evaluation-annihilator}
\end{equation}
只需证明评价映射 $\xi\mapsto\xi(z_0)$ 的像满射到 $T_pM$，即可得到
$Y_p=0$；这一点现在是显然的。于是命题得证。

下面只剩证明引理 10.5.4。

\begin{proof}[引理 10.5.4 的证明]
首要目标是证明
\begin{equation}
 \langle D_u\bar\partial_{(j,J)}\xi,\eta\rangle=0
 \label{eq:global-distributional-adjoint-zero}
\end{equation}
对所有光滑 $\xi\in\Omega^0(u^*TM)$ 成立；也就是说，$\eta$ 是
$(D_u\bar\partial_{(j,J)})^\dagger\eta=0$ \emph{在整个 $\Sigma$ 上的}分布解，而不只是在
$\Sigma\setminus\{z_0\}$ 上。

从\eqref{eq:distribution-adjoint-with-delta}出发：
\begin{equation}
 \langle D_u\bar\partial_{(j,J)}\xi,\eta\rangle
 +\langle\xi,\delta_{z_0}Y_p\rangle=0
 \quad\text{对所有 }\xi\in C^\infty.
 \label{eq:distribution-adjoint-start}
\end{equation}
已知 $\operatorname{supp}\eta\subset\{z_0\}$，先化简配对
$\langle D_u\bar\partial_{(j,J)}\xi,\eta\rangle$。
取以固定标记点 $z_0$ 为中心的复坐标 $z$，并取 $M$ 上的复坐标
$(w_1,\ldots,w_n)$，把它看作 $p$ 的某邻域上的坐标。在 $z_0$ 的邻域
$U$ 上取标准度量
\[
 h=\frac{\sqrt{-1}}2\,dz\,d\bar z,
\]
并在坐标 $(w_1,\ldots,w_n)$ 下固定 $\mathbb C^n$ 上任意埃尔米特度量。

下面的引理是证明的关键。

\begin{lemma}\label{lem:local-standard-part-pairing}
设 $\eta$ 如上。对 $u^*(TM)$ 的任意光滑截面 $\xi$ 及
$\bigl(\Omega_{1,p}^{(0,1)}(u^*TM)\bigr)^*$ 中的 $\eta$，有
\[
 \langle D_u\bar\partial_{(j,J)}\xi,\eta\rangle
 =\langle\bar\partial\xi,\eta\rangle,
\]
其中 $\bar\partial$ 是上述坐标中 $\mathbb C^n$ 上的标准柯西--黎曼算子。
\end{lemma}

\begin{proof}
已经证明 $\eta$ 是满足
$\operatorname{supp}\eta\subset\{z_0\}$ 的分布。由点支集分布结构定理 10.5.1，
\[
 \eta=P\left(\frac\partial{\partial s},
 \frac\partial{\partial t}\right)(\delta_{z_0}),
\]
其中 $z=s+it$ 是 $z_0$ 处的复坐标，而
$P(\partial/\partial s,\partial/\partial t)$ 是与二元多项式 $P$ 相伴的微分算子，其系数属于
$\bigl(\Lambda_{(j_{z_0},J_p)}^{(0,1)}(u^*TM)\bigr)^*$。
又因为 $\eta\in(W^{1,p})^*$，\emph{$P$ 的次数必须为零}，故
\begin{equation}
 \eta=\beta_{z_0}\cdot\delta_{z_0}
 \label{eq:eta-delta-mass}
\end{equation}
对某个常向量
$\beta_{z_0}\in\Lambda^{(0,1)}(T_pM)$ 成立。这是因为对
$W^{1,p}$ 映射作“某点处导数的评价”并不定义 $W^{1,p}$ 上的连续泛函。

在 $z_0$ 附近的坐标中，可写
\[
 D_u\bar\partial_{(j,J)}\xi
 =\bar\partial\xi+E\cdot\partial\xi+F\cdot\xi,
\]
其中 $E,F$ 是零阶矩阵算子，并且 $E(z_0)=0=F(z_0)$。由
\eqref{eq:eta-delta-mass}，
\[
 \langle E\cdot\partial\xi+F\cdot\xi,\eta\rangle
 =\langle E\cdot\partial\xi+F\cdot\xi,
 \beta_{z_0}\delta_{z_0}\rangle
 =\bigl(E(z_0)\partial\xi(z_0)+F(z_0)\xi(z_0),\beta_{z_0}\bigr)_{z_0}=0.
\]
所以
\[
 \langle D_u\bar\partial_{(j,J)}\xi,\eta\rangle
 =\langle\bar\partial\xi+E\cdot\partial\xi+F\cdot\xi,\eta\rangle
 =\langle\bar\partial\xi,\eta\rangle,
\]
引理得证。
\end{proof}

由这个引理，\eqref{eq:distribution-adjoint-start}化为
\begin{equation}
 \langle\bar\partial\xi,\eta\rangle
 +\langle\xi,\delta_{z_0}Y_p\rangle=0
 \quad\text{对所有 }\xi.
 \label{eq:standard-distribution-identity}
\end{equation}
把 $\xi$ 分解为
\[
 \xi(z)=\bigl(\xi(z)-\chi(z)\xi(z_0)\bigr)+\chi(z)\xi(z_0),
\]
其中 $\chi$ 是 $U$ 上的截断函数，在小邻域 $V\subset U$ 上满足
$\chi\equiv1$，且 $\operatorname{supp}\chi\subset U$。第一项
$\widetilde\xi$ 定义为
$\widetilde\xi(z):=\xi(z)-\chi(z)\xi(z_0)$，它是 $\Sigma$ 上的光滑截面，并满足
\[
 \widetilde\xi(z_0)=0,
 \qquad \bar\partial\widetilde\xi=\bar\partial\xi
 \quad\text{于 }V,
\]
因为在 $V$ 上 $\chi(z)\xi(z_0)\equiv\xi(z_0)$。在
\eqref{eq:standard-distribution-identity}中以 $\widetilde\xi$ 代替 $\xi$，得到
\[
 \langle\bar\partial\widetilde\xi,\eta\rangle
 +\langle\widetilde\xi,\delta_{z_0}Y_p\rangle=0.
\]
再利用 $\operatorname{supp}\eta\subset\{z_0\}$ 和
\eqref{eq:distribution-adjoint-start}，可得
\begin{equation}
 \langle\widetilde\xi,\delta_{z_0}Y_p\rangle
 =\langle\widetilde\xi(z_0),Y_p\rangle=0,
 \label{eq:cutoff-evaluation-vanishes}
\end{equation}
从而 $\langle\bar\partial\widetilde\xi,\eta\rangle=0$。另一方面，
\begin{equation}
 \langle\bar\partial\xi,\eta\rangle
 =\langle\bar\partial\widetilde\xi,\eta\rangle,
 \label{eq:cutoff-pairings-equal}
\end{equation}
因为在 $V$ 上 $\bar\partial\widetilde\xi=\bar\partial\xi$，而
$\operatorname{supp}\eta\subset\{z_0\}$。因此
$\langle\bar\partial\xi,\eta\rangle=0$，再由引理 10.5.5 得到
\eqref{eq:global-distributional-adjoint-zero}。

由椭圆正则性，$\eta$ 必须光滑。先前已证明
$\eta=0$ 于 $\Sigma\setminus\{z_0\}$，故由连续性，
$\eta$ 在整个 $\Sigma$ 上都为零。
\end{proof}

这个引理完成了引理 10.5.4 的证明，继而引理 10.5.4 完成命题 10.5.3 的证明。
\end{proof}

评价横截性还有另一个在粘合问题中扮演重要角色的版本，即所谓对角横截性。下面说明这个横截性命题。其证明细节只是命题 10.5.3 的变式，所以从略。
\stfhSubIndex{I289}{heng^jie^xing}{横截性（transversality）}{dui^jiao^heng^jie^xing}{对角横截性（diagonal transversality）}

回忆第 9 章：每个稳定映射极限都带有一个对偶图，用以编码稳定映射的相交型。对偶图的每个顶点对应一个不可约分支，从该顶点发出的每条边对应定义域的一个节点，而稳定映射的定义域是预稳定曲线。
由此考虑 $k$ 标记模空间
\[
 \widetilde{\mathcal M}_{g,k}(J;\beta).
\]
另取 $k$ 个模空间
$\widetilde{\mathcal M}_{g_i,\ell_i+1}(J;\alpha_i)$，每个都有一个特定标记点
$(i=1,\ldots,k)$，并以 $\operatorname{ev}_0$ 表示在该特定标记点的评价映射。
给定对偶图顶点附近的局部构形由纤维积
\begin{equation}
 \widetilde{\mathcal M}_{g,k}(J;\beta)_{\operatorname{ev}}
 \mathop{\times}_{\prod_{i=1}^k\operatorname{ev}_0^{(i)}}
 \prod_{i=1}^k\widetilde{\mathcal M}_{g_i,\ell_i+1}(J;\alpha_i)
 \label{eq:diagonal-fiber-product}
\end{equation}
给出。我们希望证明这个纤维积具有流形结构。为此考虑评价映射
\[
 \operatorname{Ev}:\widetilde{\mathcal M}_{g,k}(J;\beta)
 \times\prod_{i=1}^k\widetilde{\mathcal M}_{g_i,\ell_i+1}(J;\alpha_i)
 \longrightarrow(M\times M)^k,
\]
定义为
\begin{equation}
 \operatorname{Ev}\bigl((u,\boldsymbol z);
 \{(u_i,\boldsymbol z^{(i)})\}\bigr)
 =\prod_{i=1}^k
 \bigl(\operatorname{ev}_i(z_i),
 \operatorname{ev}_0^{(i)}(z_0^{(i)})\bigr).
 \label{eq:multi-evaluation-map}
\end{equation}

\begin{corollary}\label{cor:diagonal-fiber-product-smooth}
设所有
$\widetilde{\mathcal M}_{g,k}(J;\beta)$ 以及
$\widetilde{\mathcal M}_{g_i,\ell_i+1}(J;\alpha_i)$（$i=1,\ldots,k$）都弗雷德霍姆正则。若
$\operatorname{Ev}$ 横截于多重对角线
$\Delta_{M\times M}^k\subset(M\times M)^k$，则纤维积
\eqref{eq:diagonal-fiber-product} 光滑。
\end{corollary}
\stfhSubIndexRangeEnd{I290}{heng^jie^xing}{横截性（transversality）}{ping^jia^heng^jie^xing}{评价横截性（evaluation transversality）}

现在研究当 $J$ 通有选取时 $\operatorname{Ev}$ 的横截性。为此考虑映射
\[
 \Upsilon_k:\mathcal J_\omega\times\mathcal F_{g,k}(M;\beta)
 \times\prod_{i=1}^k\mathcal F_{g_i,\ell_i+1}(M;\alpha_i)
 \longrightarrow
 \mathcal H_g''(M;\beta)\times
 \prod_{i=1}^k\mathcal H_{g_i}''(M;\alpha_i)\times(M\times M)^k,
\]
定义为
\begin{equation}
 \Upsilon_k\bigl(J,(u,\boldsymbol z);
 \{(u_i,\boldsymbol z^{(i)})\}\bigr)
 =\bigl(\bar\partial_{(j,J)}u,
 \{\bar\partial_{(j_i,J)}u_i\}_{i=1,\ldots,k};
 \operatorname{Ev}((u,\boldsymbol z);
 \{(u_i,\boldsymbol z^{(i)})\})\bigr).
 \label{eq:universal-multi-evaluation-map}
\end{equation}
由于命题 10.5.3 的证明在每个 $z_i$ 处都是局部的，而且扰动 $B$ 可以分别局部化到每个 $z_i$，立即得到如下横截性结论。

\begin{theorem}\label{thm:universal-diagonal-transversality}
映射 $\Upsilon_k$ 横截于
\[
 o_{\mathcal H_g''(M;\beta)}\times
 \prod_{i=1}^k o_{\mathcal H_{g_i}''(M;\alpha_i)}
 \times(\Delta_M)^k
\]
在
\[
 \widetilde{\mathcal M}_{g,k}(M;\beta)\times
 \prod_{i=1}^k\widetilde{\mathcal M}_{g_i,\ell_i+1}(M;\alpha_i)
\]
上的限制。因此，
\[
 \Upsilon_k^{-1}\left(
 o_{\mathcal H_g''(M;\beta)}\times
 \prod_{i=1}^k o_{\mathcal H_{g_i}''(M;\alpha_i)}
 \times(\Delta_M)^k\right)
 \cap
 \left(\widetilde{\mathcal M}_{g,k}(M;\beta)\times
 \prod_{i=1}^k\widetilde{\mathcal M}_{g_i,\ell_i+1}(M;\alpha_i)\right)
\]
是
$\widetilde{\mathcal M}_{g,k}(M;\beta)\times
\prod_{i=1}^k\widetilde{\mathcal M}_{g_i,\ell_i+1}(M;\alpha_i)$ 的光滑子流形。此外，投影映射
\[
 \Upsilon_k^{-1}\left(
 o_{\mathcal H_g''(M;\beta)}\times
 \prod_{i=1}^k o_{\mathcal H_{g_i}''(M;\alpha_i)}
 \times(\Delta_M)^k\right)
 \cap
 \left(\widetilde{\mathcal M}_{g,k}(M;\beta)\times
 \prod_{i=1}^k\widetilde{\mathcal M}_{g_i,\ell_i+1}(M;\alpha_i)\right)
 \longrightarrow\mathcal J_\omega
\]
是指标为
\begin{equation}
 \iota(g;\beta)+2k+
 \sum_{i=1}^k\bigl(\iota(g_i;\alpha_i)+2\ell_i\bigr)
 -k\dim M
 \label{eq:diagonal-transversality-index}
\end{equation}
的弗雷德霍姆映射。
\end{theorem}

萨德--斯梅尔定理立即给出下述推论。

\begin{corollary}\label{cor:generic-diagonal-fiber-products}
存在贝尔集 $\mathcal J_\omega^{\mathrm{fiber}}\subset\mathcal J_\omega$，使得对任意
$\beta,\{\beta_i\}$ 及任意
$J\in\mathcal J_\omega^{\mathrm{fiber}}$，纤维积
\eqref{eq:diagonal-fiber-product} 与
\[
 \widetilde{\mathcal M}_{g,k}(J;\beta)\times
 \prod_{i=1}^k\widetilde{\mathcal M}_{g_i,\ell_i+1}(J;\alpha_i)
\]
的交是维数由\eqref{eq:diagonal-transversality-index}给出的光滑流形。
\end{corollary}
